\origpage{281}{GRAPH MINORS. VIII}
\noindent $G$ corresponding to $\Gamma_i$, and let $Y_i$ be the set of vertices of $G$ represented by the points in $V(\Gamma) \cap H_i$. Then $(G_i, Y_{i-1}, Y_i)$ ($1 \le i \le n+1$) is a path-decomposition of $(G, T)$ over the class of all tubular grafts with width $\ge r + 1$ and index $\le k + 2r$. It remains to check that

\hspace*{12pt}(i) for $1 \le i \le n$, $|Y_i| = r$; this is true by (ii) above and the fact that $(\Gamma, F_1, F_2)$ has width $r$

\hspace*{12pt}(ii) if $n \ge 1$, there are $r$ mutually vertex-disjoint paths of $\Gamma$ between $Y_1$ and $Y_n$; this follows from Menger's theorem and the fact that $(\Gamma, F_1, F_2)$ has width $r$.

The result follows. $\square$

\fix{We now prove (15.1) when an end has fewer than $r$ vertices. Let $A = F_1 \cap V(\Gamma)$ and $B = F_2 \cap V(\Gamma)$. Exchanging $F_1$ and $F_2$ does not change the width, so we may assume that $|A| \le |B|$, and then $|A| < r$. We identify $G$ with the graph of $\Gamma$. From now on a \textit{hoop} is a $\Gamma$-normal hoop in $\Sigma(F_1, F_2)$, and its \textit{vertices} are those of $\Gamma$ on it. A hoop or point in $\Sigma(F, F')$ lies \textit{between} $F$ and $F'$, and $F'$ comes \textit{after} $F$ if it lies between $F$ and $F_2$. A \textit{$C$-hoop} is a hoop with vertex set $C$. Since the width is $r$, and not infinite, every hoop that is neither an $A$-hoop nor a $B$-hoop has at least $r$ vertices, and some hoop $H$ is neither. So every hoop has at least $|A|$ vertices. If $F, F'$ are laminar hoops of a cylinder, every hoop in $\Sigma(F, F')$ contains $F \cap F'$, since $F$ and $F'$ lie in its two closed sides, which meet only in this hoop. We paste a disc onto each cuff of $\Sigma$ to obtain a sphere $\Sigma'$, in which each hoop $F$ bounds two closed discs, $D_1(F)$ containing the cuff beyond $F_1$ and $D_2(F)$ containing the other cuff. If $F'$ comes after $F$, then $\Sigma(F, F')$ is the set of points in the interior of neither $D_1(F)$ nor $D_2(F')$.

\hspace*{12pt}(1) \textit{Let $F, F'$ be hoops and $q \ge 0$, such that every hoop between $F$ and $F'$ has at least $q$ vertices. Then there are $q$ mutually vertex-disjoint paths of $\Gamma$ between the vertices of $F$ and those of $F'$.}

This is the use of Menger's theorem in step (ii) of the printed proof. We may assume that $F'$ comes after $F$. If there are no such paths, then by Menger's theorem some set $Z$ of fewer than $q$ vertices meets every path between the vertices of $F$ and those of $F'$ formed by vertices and edges between $F$ and $F'$. We add to $D_1(F)$ thin closed discs around the vertices joined to a vertex of $F$ by such paths that miss $Z$, and around the edges between $F$ and $F'$ at these vertices, and then the components of the rest of $\Sigma'$ that miss $D_2(F')$. These vertices are not on $F'$, and the other ends of these edges are such vertices or are in $Z$. So the boundary of the closed disc obtained is a hoop between $F$ and $F'$ whose vertices are in $Z$, a contradiction.

\hspace*{12pt}(2) \textit{Let $D$, $D'$ be closed discs in a sphere with disjoint interiors, and let $Y$ be a connected set disjoint from $D \cup D'$. Then at most two points of $bd(D) \cap bd(D')$ are in the closure of $Y$.}

Suppose that $a, b, c$ are three such points. Let $\ell_1 \subseteq D$ and $\ell_2 \subseteq D'$ be $[0,1]$-arcs from $a$ to $b$ that meet $bd(D)$ and $bd(D')$ only in $a$ and $b$. Since the interiors of $D$ and $D'$ are disjoint, $\Theta = \ell_1 \cup \ell_2$ is an $O$-arc. It divides the sphere into two open discs, and $Y$ lies in one of them, $P$, since $Y$ is connected and misses $\Theta$. The arc $\ell_1$ divides $D$ into two closed discs, and the one with $c$ in its boundary, $X_1$, satisfies $X_1 - \ell_1 \subseteq P$. For $X_1 - \ell_1$ is connected, misses $\Theta$, and contains $c$, which is in the closure of $P$ but not in $\Theta$. We define $X_2 \subseteq D'$ in the same way. Arcs from $c$ through the interiors of $X_1$ and $X_2$ to inner points $t_1$ of $\ell_1$ and $t_2$ of $\ell_2$ form a $[0,1]$-arc $\Lambda$, which divides $P$ into two open discs. Their closures meet $\Theta$ in the two arcs of $\Theta$ between $t_1$ and $t_2$, one of which contains $a$ and the other $b$. But $Y$ misses $\Lambda \subseteq D \cup D'$, so it lies in one of these discs, whose closure then contains $a$ and $b$, a contradiction.

\hspace*{12pt}(3) \textit{Let $F, F'$ be hoops such that every hoop between them has at least $r$ vertices. Then the subgraph of $G$ formed by the vertices and edges between $F$ and $F'$ has a path-decomposition $(G_i, Y_{i-1}, Y_i)$ $(1 \le i \le m)$ over the class of all tubular grafts of width $\ge r + 1$, from the vertex set of $F$ to that of $F'$, whose other sets $Y_i$ are vertex sets of hoops between $F$ and $F'$ with $r$ vertices. Each $(G_i, Y_{i-1} \cup Y_i)$ has index at most $|Y_{i-1}| + |Y_i|$.}

This is the printed argument for the tube formed by $F$, $F'$ and the part of $\Gamma$ between them, since it has width at least $r$ and its ends have at least $r$ vertices. (If its width is more than $r$, we take $m = 1$.)

\hspace*{12pt}(4) \textit{Let $F, F'$ be hoops, let $C$ be the vertex set of $F$, and suppose that every hoop between $F$ and $F'$ has at least $|C|$ vertices and some vertex $x$ of $F'$ is not in $F$. Then there are a $C$-hoop $L'$ between $F$ and $F'$, a vertex $v \notin L'$, and a hoop $L$ between $L'$ and $F'$ with vertex set $C \cup \{v\}$, such that every point of $\Gamma$ between $L'$ and $F'$ lies between $L$ and $F'$, and no hoop between $L$ and $F'$ is a $C$-hoop.}

Exchanging the names of $F_1$ and $F_2$ if necessary, we may assume that $F'$ comes after $F$. Since $F$ is a $C$-hoop, we may choose a $C$-hoop $L'$ between $F$ and $F'$ such that the number of vertices and edges of $\Gamma$ between $F$ and $L'$ is as large as possible.

First, if $F''$ is a $C$-hoop between $L'$ and $F'$, every point of $\Gamma$ between $L'$ and $F'$ lies between $F''$ and $F'$. For let $y$ be a point of $\Gamma$ between $L'$ and $F''$. The vertices and edges between $F$ and $F''$ include those between $F$ and $L'$, and so they are the same, by the choice of $L'$. If $y \notin F''$, the vertex or edge that contains $y$ misses $F''$, so it lies between $F$ and $F''$, and hence between $F$ and $L'$. Then $y \in L'$, so $y$ is a vertex in $C$, and $y \in F''$ after all.

Second, let $\Delta$ be a closed disc with $D_1(L') \subseteq \Delta \subseteq D_1(F')$ whose boundary meets $U(\Gamma)$ only in $L'$. Its boundary is a hoop between $L'$ and $F'$ with at least $|C|$ vertices, all in $C$, so it is a $C$-hoop. By the first step, every point of $\Gamma$ between $L'$ and $F'$ lies between $bd(\Delta)$ and $F'$, and so not in the interior of $\Delta$. The other points of $\Gamma$ in $\Delta$ are in $D_1(L')$.

Third, let $E_1 = D_1(L')$ and $E_2 = D_2(F')$. We find a connected set $W$ that misses $U(\Gamma) \cup E_2$ and meets the interior of $E_1$, with a vertex $v \notin E_1$ in its closure. Since $x \notin C$, the vertex $x$ of $F'$ is not on $L'$, so it is not in $E_1$. So a short $[0,1]$-arc from $x$ into $\Sigma' - E_2$ misses $E_1$, and $x$ is in the closure of the component $\Omega$ of $\Sigma' - (E_1 \cup E_2)$ that contains its other points. Some point $p \in bd(E_1)$ in the closure of $\Omega$ is not in $U(\Gamma) \cup E_2$. For otherwise $\Omega$ would be open and closed in the connected set $\Sigma' - (E_2 \cup C)$, since $\Gamma$ meets $bd(E_1)$ only in $C$, and so would contain the interior of $E_1$. Let $O$ be a connected open neighbourhood of $p$ that misses $U(\Gamma) \cup E_2$, let $M = (\Sigma' - (E_1 \cup E_2)) \cup O$, and let $W$ be the component of $M - U(\Gamma)$ that contains $O$. It misses $U(\Gamma) \cup E_2$, and it meets the interior of $E_1$, since $O$ does. Suppose that no vertex outside $E_1$ is in the closure of $W$. Then every edge $e$ of $\Gamma$ in the closure of $W$ that misses $E_1 \cup E_2$ has its ends in $E_1$, by our assumption, and so on $bd(E_1)$ and in $C$. If $e$ has two ends, its closure divides $\Sigma' - (E_1 - bd(E_1))$ into two closed discs, and adding to $E_1$ the one that misses the interior of $E_2$ gives a closed disc with $e$ in its boundary. Pushing this boundary a little beyond $e$, off $U(\Gamma) \cup E_2$, gives a disc $\Delta$ as in the second step with $e$ in its interior, which is impossible. So $e$ is a loop, and by the same argument the closed disc $D_e$ bounded by the closure of $e$ that misses the interior of $E_1$ also misses the interior of $E_2$. It meets $E_1 \cup E_2$ only in the end of $e$, since the other points of $E_1 \cup E_2$ are in the closure of the interiors of $E_1$ and $E_2$ and not in the closure of $e$. Let $N$ be the union of $W$ and of the sets $D_e$ less the ends of $e$. Near a point of such an edge $e$, one side of $e$ is in $W$ and the other in $D_e$, so $N$ is open. Also $N$ is closed in $M$. Since the ends of the discs $D_e$ are not in $M$, it is enough to show that each $y \in M - W$ in the closure of $W$ is in $N$. Since $W$ is a component of the open set $M - U(\Gamma)$, $y$ is in $U(\Gamma)$. By our assumption $y$ is not a vertex, so it lies on an edge that misses $E_1 \cup E_2$. This edge is in the closure of $W$, again since $W$ is a component of $M - U(\Gamma)$, so $y \in N$. The set $\Omega \cup O \subseteq M$ is connected, since $O$ meets $\Omega$, and it meets $N$, so it lies in $N$. Hence $x$ is in the closure of $W$ or in some $D_e$. Both are false, the second since $x \in E_2$ and $x \notin C$. So some vertex $v \notin E_1$ is in the closure of $W$.

Fourth, $v$ lies between $L'$ and $F'$ and not on $L'$, since $v \notin E_1$ and $W$ misses $E_2$. The region $R$ of the drawing formed by $\Gamma$ and $F'$ that contains $W$ misses $E_2$, since it misses $F'$ and meets $E_1$. Since $v$ is in the closure of $R$, there is a $[0,1]$-arc from $v$ with its other points in $R$. We continue it in $R$ to a point of $W$ in the interior of $E_1$, stopping where it first meets $E_1$, which is on $L'$. We add to $E_1$ a thin closed disc around the arc so obtained, which meets $U(\Gamma) \cup E_2$ only in $v$, and meets $E_1$ in a short arc of $L'$ that misses $U(\Gamma)$. The boundary $L$ of the union is a hoop between $L'$ and $F'$ with vertex set $C \cup \{v\}$. Every point of $\Gamma$ in the union is in $E_1$ or is $v$. So every point of $\Gamma$ between $L'$ and $F'$ either is outside the union, and so lies between $L$ and $F'$, or is in $C \cup \{v\}$, and so is on $L$.

Finally, suppose that a hoop $F''$ between $L$ and $F'$ is a $C$-hoop. It lies between $L'$ and $F'$, so by the first step $v$ lies between $F''$ and $F'$. Also $v$ lies between $L'$ and $F''$, since $v \in L$ and $L$ lies between $L'$ and $F''$. So $v \in F''$, and $v \in C$, contrary to $v \notin L'$. This proves (4).

\hspace*{12pt}(5) \textit{Let $F, F'$ be hoops with the same vertex set $C$, and let $J \supseteq C$ be a set of vertices, none of which except those in $C$ lies between $F$ and $F'$. Then the subgraph of $G$ formed by $J$ and the vertices and edges between $F$ and $F'$, with root set $J$, has a path-decomposition over the class of all tubular grafts of width $\ge |J| + 1$ and index $|J|$, with all $Y_i$ equal to $J$. This class contains the graft with vertex set $J$, no edges and root set $J$.}

We may assume that $F'$ comes after $F$. Since $\Gamma$ meets $F \cup F'$ only in $C$, each component $Y$ of the set of points of $U(\Gamma) - C$ between $F$ and $F'$ misses $D_1(F) \cup D_2(F')$. These discs have disjoint interiors, and both boundaries contain $C$. So by (2), at most two vertices in $C$, the \textit{attachments} of $Y$, are in the closure of $Y$. Let $G_Y$ be the subgraph formed by $J$ and the vertices and edges of $\Gamma$ in the closure of $Y$, and let $G_0$ be the graph with vertex set $J$ and no edges. Every edge between $F$ and $F'$ is in exactly one $G_Y$, and any two of these graphs meet exactly in $J$. So $(G_0, J, J)$, followed by the triples $(G_Y, J, J)$, is the required path-decomposition, if $(G_0, J)$ and each $(G_Y, J)$ are tubular of width $\ge |J| + 1$.

To show this, let $P$ be the closure of $Y$, or the null set for $G_0$. It lies in $D_2(F)$ and meets $F$ only in attachments of $Y$. Choose distinct points $p, q \in F$ with these attachments among them, and a $[0,1]$-arc $\lambda$ from $p$ to $q$ that meets $D_2(F)$ only in $p$ and $q$. It divides $D_1(F)$ into two closed discs $K_1$, $K_2$. We draw the vertices of $J$ not in $P$ at distinct points of $\lambda - \{p, q\}$, and remove a small open disc from the interior of each of $K_1$, $K_2$. In the cylinder so obtained, $bd(K_1)$ and $bd(K_2)$ are laminar hoops, and with this drawing they form a tube. It frames $(G_Y, J)$, or $(G_0, J)$, since the drawing meets $bd(K_1) \cup bd(K_2)$ exactly in $J$. By the fact above, every hoop between $bd(K_1)$ and $bd(K_2)$ contains $bd(K_1) \cap bd(K_2) = \lambda$, so its vertex set is $J$ or has more than $|J|$ elements. So this tube has width at least $|J| + 1$. This proves (5).

We now complete the proof. Suppose first that $|B| \ge r$. Since $|B| > |A|$, some vertex of $F_2$ is not in $F_1$, so (4) applies with $F = F_1$ and $F' = F_2$. Let $L_1'$, $v_1$ and $L_1$ be as given there, and let $L = L_1$, $J = A \cup \{v_1\}$ and $R = F_2$. The hoop $L$ is not an $A$-hoop and has $|A| + 1 \le r$ vertices. If it is a $B$-hoop, then $|A| + 1 = |B| \ge r$, and otherwise it has at least $r$ vertices. So $|J| = r$. Every hoop between $L$ and $R$ has at least $r$ vertices, since it is not an $A$-hoop, by (4), and $|B| \ge r$.

Suppose now that $|B| < r$, so that $H$ has at least $r$ vertices. We choose a hoop $H_0$ such that every hoop between $H_0$ and $F_2$ has at least $|B|$ vertices, some vertex of $H_0$ is not in $F_2$, and $H_0 = F_1$, or $A = B$, or $|A| + 1 < r$. If $A = B$, let $H_0 = H$. It contains $F_1 \cap F_2$, and so it contains $A$ and has a vertex not in $A = B$. If $A \ne B$ and $|A| = |B|$, let $H_0 = F_1$, which has a vertex not in $B$, since $A \not\subseteq B$. In these two cases every hoop has at least $|A| = |B|$ vertices. If $|A| < |B|$, then $|A| + 1 < r$, and we let $H_0 = H$. We claim that $B \subseteq H$. Then every hoop between $H$ and $F_2$ contains $H \cap F_2$, and so contains $B$, and $H$ has a vertex not in $B$, since it is not a $B$-hoop. Suppose that some $y \in B$ is not on $H$. Since $H$ has more than $|A|$ vertices, (4) applies with $F = F_1$ and $F' = H$, and gives a hoop between $F_1$ and $H$ with vertex set $A \cup \{v\}$, where $v \notin A$. This hoop is not an $A$-hoop and has $|A| + 1 < r$ vertices, so it is a $B$-hoop, and $y$ is on it. So $y$ lies between $F_1$ and $H$, and also between $H$ and $F_2$, since $y \in F_2$. Hence $y \in H$, a contradiction.

Next, (4) applies with $F = F_2$, $F' = H_0$ and $C = B$. It gives a $B$-hoop $R'$, a vertex $w \notin R'$ and a hoop $R$ between $H_0$ and $R'$ with vertex set $K = B \cup \{w\}$, such that every point of $\Gamma$ between $H_0$ and $R'$ lies between $H_0$ and $R$, and no hoop between $H_0$ and $R$ is a $B$-hoop. The hoop $R$ is not a $B$-hoop, since $w \notin B$. It is not an $A$-hoop either, for otherwise $B \subseteq A$, so $A = B$, since $|A| \le |B|$, and $w \in B$. So $r \le |K| = |B| + 1 \le r$.

Next, (4) applies with $F = F_1$ and $F' = R$, since $|K| > |A|$, and gives $L_1'$, $v_1$, and $L_1$ with vertex set $A \cup \{v_1\}$. If $|A| + 1 = r$, let $L = L_1$ and $J = A \cup \{v_1\}$. A hoop between $L$ and $R$ is not an $A$-hoop, by (4). It is not a $B$-hoop either, by the choice of $R$ if $H_0 = F_1$, and since it is not an $A$-hoop if $A = B$. The third choice, $|A| + 1 < r$, does not hold. So this hoop has at least $r$ vertices. If $|A| + 1 < r$, then $L_1$ is not an $A$-hoop and has fewer than $r$ vertices, so it is a $B$-hoop. The hoops between $L_1$ and $R$ are not $A$-hoops, by (4), so they have at least $|B|$ vertices, and $w$ is a vertex of $R$ not on $L_1$. So (4) applies with $F = L_1$, $F' = R$ and $C = B$, and gives $L_2'$, $v_2$, and $L_2$ with vertex set $J = B \cup \{v_2\}$. Let $L = L_2$. Then $|J| = |B| + 1 = r$, and a hoop between $L$ and $R$ is neither an $A$-hoop nor a $B$-hoop, so it has at least $r$ vertices.

In all cases $|J| = r$, and every hoop between $L$ and $R$ has at least $r$ vertices. The part of $G$ between $F_1$ and $L$ has the path-decomposition formed by $(G_0, A, J)$, where $G_0$ has vertex set $J$ and no edges, followed by that of (5) for $F_1$, $L_1'$ and $J$, and, if $L = L_2$, by that of (5) for $L_1$, $L_2'$ and $J$. Here (5) applies, since $v_1$ lies between $L_1'$ and $R$ and not on $L_1'$, and $v_2$ lies between $L_2'$ and $R$ and not on $L_2'$. These graphs contain all edges between $F_1$ and $L$, since by (4) the points of $\Gamma$ between $L_1'$ and $L_1$ also lie between $L_1$ and $R$, so they are on $L_1$, and in the same way those between $L_2'$ and $L_2$ are on $L_2$. If $R \ne F_2$, the part between $R$ and $F_2$ has the path-decomposition formed by that of (5) for $F_2$, $R'$ and $K$, followed by $(G_0', K, B)$, where $G_0'$ has vertex set $K$ and no edges. Here (5) applies, since $w$ lies between $H_0$ and $R'$ and not on $R'$, and these graphs contain all edges between $R$ and $F_2$, since by (4) the points of $\Gamma$ between $R$ and $R'$ also lie between $H_0$ and $R$, so they are on $R$. The part between $L$ and $R$ has that of (3). The graphs of one part meet those of another only in vertices of $L$ or $R$, and $T = A \cup B$. So these path-decompositions, in order, form one of $(G, T)$. Its pieces have width $\ge r + 1$. Their index is $r$, except for those from (3), whose index is at most $2r + |B| \le k + 2r$, since $B \subseteq T$. Its sets $Y_i$ other than the first and the last have $r$ elements. The second is $J$, and the second to last is the vertex set of a hoop between $L$ and $R$, so by (1) there are $r$ disjoint paths between them. So $(G, T)$ is an $r$-train over the class of all tubular grafts of width $\ge r + 1$ and index $\le k + 2r$. $\square$}

\stmt
(15.2) \textit{Let $(G, T)$ be a tubular graft of index $\le k$ and let $r \ge 0$ be an integer. Then $(G, T)$ is a concatenation of tubular grafts of width $\ge r$ and index $\le k + r(r - 1)$.}

\textit{Proof.} Let $\mathcal{G}(t, u)$ be the class of all tubular grafts of width $\ge t$ and index $\le u$. Now by (15.1),
\begin{equation*}
\mathcal{G}(0, k) \subseteq L^0(\mathcal{G}(1, k)) = C^0(\mathcal{G}(1, k)).
\end{equation*}
We prove by induction on $t$ that for $t \ge 0$,
\begin{equation*}
\mathcal{G}(0, k) \subseteq C^t(\mathcal{G}(t + 1, k + (t + 1)t)).
\end{equation*}
The result is true if $t = 0$, and we assume $t \ge 1$. By our inductive hypothesis,
\begin{equation*}
\mathcal{G}(0, k) \subseteq C^{t-1}(\mathcal{G}(t, k + t(t - 1)))
\end{equation*}
but by (15.1),
\begin{equation*}
\mathcal{G}(t, k + t(t - 1)) \subseteq L^t(\mathcal{G}(t + 1, k + (t + 1)t)),
\end{equation*}
and so
\begin{equation*}
\mathcal{G}(0, k) \subseteq C^{t-1}(L^t(\mathcal{G}(t + 1, k + (t + 1)t))) = C^t(\mathcal{G}(t + 1, k + (t + 1)t)).
\end{equation*}
This completes the inductive proof, and the result follows. $\square$

Let $\mathcal{A}(k, r, s)$ denote the class of all tubular grafts with index $\le k$ which have a framing $(\Gamma, F_1, F_2)$ in a cylinder $\Sigma$, such that there do not exist $\Gamma$-normal laminar hoops $F_1', F_2' \subseteq \Sigma(F_1, F_2)$ such that $(\Gamma', F_1', F_2')$ has width $\ge r$ and length $\ge s$, where $\Gamma'$ is the intersection of $\Gamma$ and $\Sigma(F_1', F_2')$. Our next object is to prove that $\mathcal{A}(k, r, s)$ is well-rooted.

\stmt
(15.3) \textit{If $(G, T) \in \mathcal{A}(k, r, s)$, then $(G, T)$ is a concatenation of tubular grafts of length $< s$ and index $\le k + r(r - 1)$.}
